\documentclass[11pt]{article}

\usepackage[a4paper,margin=31mm]{geometry}
\usepackage{amsmath,amssymb,amsthm,mathtools}
\usepackage{microtype}
\usepackage[hidelinks]{hyperref}
\hypersetup{
  pdftitle={The sharp lower-limit constant for the first decrease of a harmonic denominator},
  pdfauthor={Wouter van Doorn}
}

\newtheorem{theorem}{Theorem}
\newtheorem{lemma}[theorem]{Lemma}
\newtheorem{corollary}[theorem]{Corollary}

\newcommand{\Z}{\mathbb Z}
\newcommand{\F}{\mathbb F}
\newcommand{\lcm}{\mathop{\rm lcm}}

\title{The sharp lower-limit constant for the first decrease\\
of a harmonic denominator}
\author{ChatGPT}
\date{}

\begin{document}
\maketitle

\begin{abstract}
For positive integers $a\leq b$, write
\[
  \sum_{j=a}^{b}\frac1j=\frac{u_{a,b}}{v_{a,b}},
  \qquad \gcd(u_{a,b},v_{a,b})=1,
  \qquad v_{a,b}>0,
\]
and let $b(a)$ be the least $b>a$ for which $v_{a,b}<v_{a,b-1}$.
For
\[
  f_d(X)=\sum_{i=0}^{d}\prod_{\substack{0\leq j\leq d\\j\neq i}}(X-j)
\]
and for the density $\delta_d$ of primes modulo which $f_d$ has a root, put
\[
  c=\sum_{d=1}^{\infty}\frac{\delta_d}{d(d+1)}.
\]
We prove that, for every $\varepsilon>0$, there are infinitely many pairs $a<b$ with
\[
  b<a+\left(\frac1{1+c}+\varepsilon\right)\log a
  \qquad\text{and}\qquad
  v_{a,b}<v_{a,b-1}.
\]
Together with the reverse inequality established in the accompanying manuscript~\cite{VD}, this gives
\[
  \liminf_{a\to\infty}\frac{b(a)-a}{\log a}=\frac1{1+c}.
\]
The main additional input is a finite-field form of Hal\'asz's inequality; we also use simultaneous Dirichlet approximation in an elementary auxiliary construction.
\end{abstract}

\section{Notation and the main theorem}

For $a\leq b$, define
\[
  L_{a,b}=\lcm(a,a+1,\ldots,b),
  \qquad
  X_{a,b}=L_{a,b}\sum_{j=a}^{b}\frac1j,
  \qquad
  g_{a,b}=\gcd(X_{a,b},L_{a,b}).
\]
Then
\[
  v_{a,b}=\frac{L_{a,b}}{g_{a,b}}.
\]
We use $\nu_p$ for the $p$-adic valuation. All logarithms are natural.

For $d\geq0$, let
\[
  f_d(X)=\sum_{i=0}^{d}\prod_{\substack{0\leq j\leq d\\j\neq i}}(X-j).
\]
In particular, $f_0=1$. Whenever $x\notin\{0,1,\ldots,d\}$ in a field,
\begin{equation}\label{eq:partial-fraction}
  \sum_{i=0}^{d}\frac1{x-i}
  =\frac{f_d(x)}{\prod_{i=0}^{d}(x-i)}.
\end{equation}
Let $\delta_d$ be the density of primes $p$ for which $f_d$ has a root modulo $p$. This density exists by Chebotarev's theorem~\cite{Chebotarev}. Set
\begin{equation}\label{eq:def-c}
  c=\sum_{d=1}^{\infty}\frac{\delta_d}{d(d+1)}.
\end{equation}

\begin{theorem}\label{thm:main}
For every $\varepsilon>0$, there are infinitely many positive integers $a<b$ such that
\[
  b<a+\left(\frac1{1+c}+\varepsilon\right)\log a
\]
and
\[
  v_{a,b}<v_{a,b-1}.
\]
Consequently,
\[
  \liminf_{a\to\infty}\frac{b(a)-a}{\log a}\leq\frac1{1+c}.
\]
\end{theorem}

\begin{corollary}
Combined with the lower bound from~\cite{VD},
\[
  \liminf_{a\to\infty}\frac{b(a)-a}{\log a}\geq\frac1{1+c},
\]
Theorem~\ref{thm:main} gives
\[
  \liminf_{a\to\infty}\frac{b(a)-a}{\log a}=\frac1{1+c}.
\]
\end{corollary}

\section{Local denominator calculations}

We first isolate the elementary facts about adding the endpoint $1/b$.

\begin{lemma}\label{lem:endpoint}
Let $U/V$ be a reduced fraction with $V>0$, and let $V'$ be the denominator in lowest terms of $U/V+1/b$. Then, for every prime $p$,
\[
  \nu_p(V')-\nu_p(V)\leq \nu_p(b).
\]
\end{lemma}

\begin{proof}
Put $s=\nu_p(V)$ and $k=\nu_p(b)$. Since $p\nmid U$,
\[
  \frac UV+\frac1b=\frac{bU+V}{bV}.
\]
If $k<s$, then $\nu_p(bU+V)=k$, so the new denominator exponent is $s$. If $k>s$, then $\nu_p(bU+V)=s$, so the new denominator exponent is $k$. If $k=s$, the new denominator exponent is at most $s$. The assertion follows in all three cases.
\end{proof}

\begin{lemma}\label{lem:poly}
For every $d\geq0$,
\begin{equation}\label{eq:symmetry}
  f_d(d-X)=(-1)^d f_d(X).
\end{equation}
Moreover, if $p>d$ is prime and $f_d(x)=0$ in $\F_p$, then
\begin{equation}\label{eq:root-avoidance}
  x\notin\{0,1,\ldots,d\}.
\end{equation}
\end{lemma}

\begin{proof}
Replacing $j$ by $d-j$ gives
\[
\begin{aligned}
  f_d(d-X)
  &=\sum_{i=0}^{d}\prod_{\substack{0\leq j\leq d\\j\neq i}}(d-X-j)\\
  &=(-1)^d\sum_{i=0}^{d}\prod_{\substack{0\leq j\leq d\\j\neq i}}(X-j)
   =(-1)^d f_d(X).
\end{aligned}
\]
If $x=i\in\{0,\ldots,d\}$, every summand in $f_d(i)$ vanishes except the summand indexed by $i$, and
\[
  f_d(i)=\prod_{\substack{0\leq j\leq d\\j\neq i}}(i-j)\neq0\pmod p
\]
because $p>d$.
\end{proof}

The following two lemmas concern primes larger than the square root of the interval length.

\begin{lemma}[A soluble prime gives a decrease]\label{lem:S-local}
Let $n=b-a$, let $\sqrt n<p\leq n$ be prime, and put $d=\lfloor n/p\rfloor$. Suppose that $p\mid b$ and, with
\[
  z\equiv \frac bp\pmod p,
\]
we have $f_d(z)=0$ in $\F_p$. Then
\[
  \nu_p(v_{a,b})\leq \nu_p(v_{a,b-1})-1.
\]
\end{lemma}

\begin{proof}
By Lemma~\ref{lem:poly}, $z\notin\{0,\ldots,d\}$. The multiples of $p$ in $[a,b]$ are
\[
  b,b-p,\ldots,b-dp.
\]
For $0\leq i\leq d$,
\[
  \frac{b-ip}{p}\equiv z-i\not\equiv0\pmod p.
\]
Thus none of these integers is divisible by $p^2$. Hence
\[
  \nu_p(L_{a,b})=\nu_p(L_{a,b-1})=1.
\]
Modulo $p$, only the displayed multiples of $p$ contribute to $X_{a,b}$, and therefore, using \eqref{eq:partial-fraction},
\[
\begin{aligned}
  X_{a,b}
  &\equiv \frac{L_{a,b}}p\sum_{i=0}^{d}\frac1{z-i}\\
  &\equiv \frac{L_{a,b}}p
       \frac{f_d(z)}{\prod_{i=0}^{d}(z-i)}
   \equiv0\pmod p.
\end{aligned}
\]
On the other hand, $L_{a,b}/b$ is a $p$-adic unit. Since
\[
  X_{a,b}
  =\frac{L_{a,b}}{L_{a,b-1}}X_{a,b-1}+\frac{L_{a,b}}b
\]
and $L_{a,b}/L_{a,b-1}$ is also a $p$-adic unit, it follows that $p\nmid X_{a,b-1}$. Thus the old denominator contains one factor $p$, while the new denominator contains none.
\end{proof}

\begin{lemma}[The bad residues for an insoluble prime]\label{lem:T-local}
Let $n=b-a$, let $\sqrt n<p\leq n$ be prime, and put $d=\lfloor n/p\rfloor$. Suppose that $p\mid b$ and that $f_d$ has no root in $\F_p$. Define
\begin{equation}\label{eq:bad-z}
  \mathcal Z_{p,d}
  =\{0\}\cup
   \bigl\{z\in\F_p\setminus\{0,1,\ldots,d\}:
                 f_{d-1}(z-1)=0\bigr\}.
\end{equation}
If
\[
  z\equiv \frac bp\pmod p
  \qquad\text{satisfies}\qquad
  z\notin\mathcal Z_{p,d},
\]
then
\[
  \nu_p(v_{a,b})\leq \nu_p(v_{a,b-1}).
\]
Moreover,
\begin{equation}\label{eq:bad-cardinality}
  |\mathcal Z_{p,d}|\leq d.
\end{equation}
\end{lemma}

\begin{proof}
First suppose that $z\in\{1,\ldots,d\}$. Then $b-zp$ is divisible by $p^2$. Since $p^2>n$, it is the unique multiple of $p^2$ in $[a,b]$, and it belongs to both $[a,b-1]$ and $[a,b]$. If $p^k\Vert b-zp$, then this is the unique term of maximal $p$-adic order in both intervals. After multiplication by the corresponding least common multiple, that term is a $p$-adic unit and every other term is divisible by $p$. Hence both $X_{a,b-1}$ and $X_{a,b}$ are nonzero modulo $p$, and the $p$-parts of the two reduced denominators are equal.

Now suppose that $z\notin\{0,1,\ldots,d\}$. Then none of
\[
  b,b-p,\ldots,b-dp
\]
is divisible by $p^2$, and
\[
  \nu_p(L_{a,b})=\nu_p(L_{a,b-1})=1.
\]
Since $f_d$ has no root modulo $p$, equation \eqref{eq:partial-fraction} gives
\[
  X_{a,b}
  \equiv \frac{L_{a,b}}p
      \frac{f_d(z)}{\prod_{i=0}^{d}(z-i)}
  \not\equiv0\pmod p.
\]
Similarly,
\[
\begin{aligned}
  X_{a,b-1}
  &\equiv \frac{L_{a,b-1}}p\sum_{i=1}^{d}\frac1{z-i}\\
  &\equiv \frac{L_{a,b-1}}p
      \frac{f_{d-1}(z-1)}{\prod_{i=1}^{d}(z-i)}
  \not\equiv0\pmod p,
\end{aligned}
\]
because $z\notin\mathcal Z_{p,d}$. Again the two $p$-parts are equal.

Finally, $f_{d-1}$ has degree $d-1$ and is not the zero polynomial modulo $p$, since its leading coefficient is $d$ and $p>d$. Hence the second set in \eqref{eq:bad-z} has at most $d-1$ elements, proving \eqref{eq:bad-cardinality}.
\end{proof}

We shall also use a fixed collection of small primes, but at growing prime-power scales.

\begin{lemma}[A prime-power-scale decrease]\label{lem:power-local}
Let $\ell>3$ be prime, let $m=\ell^e$, and suppose
\[
  2m\leq n=b-a<3m,
  \qquad m\mid b.
\]
Put $z=b/m$. If $f_2(z)=0$ in $\F_\ell$, then
\[
  \nu_\ell(v_{a,b})\leq \nu_\ell(v_{a,b-1})-1.
\]
\end{lemma}

\begin{proof}
The multiples of $m$ in $[a,b]$ are exactly
\[
  b,b-m,b-2m.
\]
By Lemma~\ref{lem:poly}, $z,z-1,z-2$ are nonzero modulo $\ell$. Thus none of these three integers is divisible by $\ell^{e+1}$, and
\[
  \nu_\ell(L_{a,b})=\nu_\ell(L_{a,b-1})=e.
\]
Modulo $\ell$,
\[
\begin{aligned}
  X_{a,b}
  &\equiv \frac{L_{a,b}}m
      \left(\frac1z+\frac1{z-1}+\frac1{z-2}\right)\\
  &\equiv \frac{L_{a,b}}m
      \frac{f_2(z)}{z(z-1)(z-2)}
  \equiv0\pmod\ell.
\end{aligned}
\]
The endpoint term $L_{a,b}/b$ is an $\ell$-adic unit. In
\[
  X_{a,b}
  =\frac{L_{a,b}}{L_{a,b-1}}X_{a,b-1}+\frac{L_{a,b}}b,
\]
the first multiplier is an $\ell$-adic unit, so $\ell\nmid X_{a,b-1}$. The asserted denominator decrease follows.
\end{proof}

\section{Prime products}

Fix a large integer $n$. For every integer $d$ with $1\leq d\leq\sqrt n-1$, let
\[
  S_d=S_d(n)
  =\left\{p:\frac{n}{d+1}<p\leq\frac nd,
                    \ f_d\text{ has a root modulo }p\right\},
\]
and let $T_d=T_d(n)$ be the complementary set of primes in the same interval. Put
\begin{equation}\label{eq:Q}
  S=\bigcup_{1\leq d\leq\sqrt n-1}S_d,
  \qquad
  Q=\prod_{q\in S}q.
\end{equation}
For a fixed positive integer $D$, put
\begin{equation}\label{eq:P}
  T^{(D)}=\bigcup_{d=1}^{D}T_d,
  \qquad
  P=\prod_{p\in T^{(D)}}p,
  \qquad
  \tau_D=\sum_{d=1}^{D}\frac{1-\delta_d}{d(d+1)}.
\end{equation}

\begin{lemma}\label{lem:products}
As $n\to\infty$,
\begin{equation}\label{eq:Q-asymp}
  \log Q=(c+o(1))n.
\end{equation}
For each fixed $D$,
\begin{equation}\label{eq:P-asymp}
  \log P=(\tau_D+o(1))n.
\end{equation}
Moreover,
\begin{equation}\label{eq:tau-limit}
  \tau_D\longrightarrow1-c
  \qquad(D\to\infty).
\end{equation}
\end{lemma}

\begin{proof}
For each fixed $d$, Chebotarev's theorem together with partial summation gives
\[
  \sum_{p\in S_d}\log p
  =\left(\frac{\delta_d}{d(d+1)}+o(1)\right)n.
\]
Fix $R$. Summing this for $d\leq R$ gives the corresponding initial part of \eqref{eq:Q-asymp}. All primes belonging to $S_d$ with $d>R$ are at most $n/(R+1)$, so the prime number theorem bounds their total logarithm by
\[
  \sum_{p\leq n/(R+1)}\log p=O\!\left(\frac nR\right).
\]
First let $n\to\infty$ and then $R\to\infty$. This proves \eqref{eq:Q-asymp}. Equation \eqref{eq:P-asymp} follows in the same way, now with only the fixed set $d\leq D$, and using the complementary density $1-\delta_d$. Finally,
\[
  \sum_{d=1}^{\infty}\frac1{d(d+1)}=1,
\]
so \eqref{eq:tau-limit} follows from \eqref{eq:def-c}.
\end{proof}

\section{Auxiliary prime powers and the choice of interval lengths}

Fix $\varepsilon>0$. By Lemma~\ref{lem:products}, we may choose $D$ so large that
\begin{equation}\label{eq:D-choice}
  \frac1{\tau_D+2c}
  <\frac1{1+c}+\frac\varepsilon2.
\end{equation}
Choose an even integer
\begin{equation}\label{eq:E-choice}
  E>D+1.
\end{equation}
Choose an integer $r\geq2$ such that
\begin{equation}\label{eq:r-choice}
  2^r>D,
\end{equation}
and choose distinct primes
\[
  \ell_1,\ldots,\ell_r\equiv1\pmod{12}.
\]
For every such prime, $3$ is a quadratic residue, and
\[
  f_2(X)=3X^2-6X+2=3(X-1)^2-1
\]
has exactly two roots modulo $\ell_i$.

\begin{lemma}\label{lem:simultaneous-powers}
There are infinitely many tuples of positive integers $(e_1,\ldots,e_r)$ such that, with
\[
  m_i=\ell_i^{e_i},
\]
we have
\begin{equation}\label{eq:powers-close}
  \frac{\max_i m_i}{\min_i m_i}\longrightarrow1
\end{equation}
as the tuples tend to infinity.
\end{lemma}

\begin{proof}
For $2\leq i\leq r$, put
\[
  \alpha_i=\frac{\log\ell_1}{\log\ell_i}.
\]
By simultaneous Dirichlet approximation, for every $H$ there is an integer $1\leq q\leq H^{r-1}$ such that
\[
  \max_{2\leq i\leq r}\|q\alpha_i\|\leq\frac1H.
\]
Along a sequence with $H\to\infty$, the corresponding $q$ are unbounded. Indeed, if they were bounded, one fixed positive integer $q$ would occur for arbitrarily large $H$, forcing $q\alpha_i\in\Z$ for every $i$. This is impossible because $\log\ell_1/\log\ell_i$ is irrational for distinct primes $\ell_1$ and $\ell_i$. Take $e_1=q$ and take $e_i$ to be the nearest integer to $q\alpha_i$. Then
\[
  e_i\log\ell_i-e_1\log\ell_1=o(1)
\]
uniformly in $i$, which is equivalent to \eqref{eq:powers-close}.
\end{proof}

For one of these tuples, put
\begin{equation}\label{eq:n-N}
  M=\max_i m_i,
  \qquad
  n=2M,
  \qquad
  N=\prod_{i=1}^{r}m_i,
  \qquad
  G=\prod_{i=1}^{r}\ell_i.
\end{equation}
For all sufficiently large tuples,
\begin{equation}\label{eq:scale-window}
  2m_i\leq n<3m_i
  \qquad(1\leq i\leq r),
\end{equation}
and
\begin{equation}\label{eq:N-small}
  \log N=O(\log n)=o(n).
\end{equation}
We now form $P=P(n)$ and $Q=Q(n)$ as in \eqref{eq:Q} and \eqref{eq:P}. Since every prime occurring in $P Q$ is larger than $\sqrt n$, while the primes $\ell_i$ are fixed, for all sufficiently large tuples the three integers $N$, $P$, and $Q$ are pairwise coprime.

We shall eventually take
\begin{equation}\label{eq:b-form}
  b=NPQY,
  \qquad
  a=b-n.
\end{equation}
For each $i$, the factor
\[
  A_i=\frac{N}{m_i}PQ
\]
is invertible modulo $\ell_i$. Requiring
\[
  A_iY\pmod{\ell_i}
\]
to be one of the two roots of $f_2$ gives two allowed nonzero residue classes for $Y$ modulo $\ell_i$. By the Chinese remainder theorem, these conditions define a set
\begin{equation}\label{eq:A-set}
  \mathcal A=\mathcal A(n)\subseteq\Z/G\Z,
  \qquad
  |\mathcal A|=2^r>D.
\end{equation}
If $Y\pmod G$ belongs to $\mathcal A$, then Lemma~\ref{lem:power-local} and \eqref{eq:scale-window} give
\begin{equation}\label{eq:aux-drop}
  \nu_{\ell_i}(v_{a,b})
  \leq\nu_{\ell_i}(v_{a,b-1})-1
  \qquad(1\leq i\leq r).
\end{equation}

\section{A bounded-width family of CRT solutions}

The set $S_E$ contains
\begin{equation}\label{eq:reservoir-size}
  K=|S_E|
  =\left(\frac{\delta_E}{E(E+1)}+o(1)\right)\frac n{\log n}
\end{equation}
primes. Here $\delta_E>0$: for example, primes splitting completely in the splitting field of $f_E$ form a positive-density subset of the primes for which $f_E$ has a root. Since $E$ is even, Lemma~\ref{lem:poly} gives
\[
  f_E(E-X)=f_E(X).
\]
For every sufficiently large $q\in S_E$, each root $x$ is therefore accompanied by the distinct root $E-x$. Distinctness follows from \eqref{eq:root-avoidance}, because equality would give $x=E/2$.

For every $q\in S_d$, we require
\begin{equation}\label{eq:S-congruence}
  NP\frac QqY\equiv x_q\pmod q,
\end{equation}
where $x_q$ is a root of $f_d$ modulo $q$. For $q\in S_E$, choose a paired set of roots $x_q^+,x_q^-$. They give two allowed residues
\[
  y_q^\pm
  \equiv x_q^\pm\left(NP\frac Qq\right)^{-1}\pmod q
\]
for $Y\pmod q$. Put, in $\F_q$,
\[
  m_q=\frac{y_q^++y_q^-}{2},
  \qquad
  h_q=\frac{y_q^+-y_q^-}{2}.
\]
Thus $h_q\neq0$.

Let $e_q$ be the CRT idempotent modulo $Q$ satisfying
\[
  e_q\equiv1\pmod q,
  \qquad
  e_q\equiv0\pmod{Q/q}.
\]
Choose the integer $C_q$ with
\begin{equation}\label{eq:Cq}
  C_q\equiv h_qe_q\pmod Q,
  \qquad
  -\frac Q2<C_q<\frac Q2.
\end{equation}
Since $Q/q\mid C_q$, there is an integer $t_q$ such that
\begin{equation}\label{eq:tq}
  C_q=\frac Qq t_q,
  \qquad
  0<|t_q|<\frac q2.
\end{equation}

Choose one prime $q_0\in S_E$. If necessary, choose one further prime of $S_E$ and fix one of its two roots, so that an even number $2J$ of reservoir primes remains. By \eqref{eq:reservoir-size},
\begin{equation}\label{eq:J-size}
  J\asymp_E\frac n{\log n}.
\end{equation}
No two of these centered integers are equal: if $C_q=C_r$ for distinct primes $q,r$, then the common value is divisible by both $Q/q$ and $Q/r$, hence by $Q$, contradicting \eqref{eq:Cq} and $h_q\neq0$. Order the corresponding centered integers $C_q$ increasingly and pair adjacent terms. If the $j$th pair is $(q_j,r_j)$, define
\begin{equation}\label{eq:Dj}
  D_j=C_{r_j}-C_{q_j}>0.
\end{equation}
Then
\begin{equation}\label{eq:width}
  \sum_{j=1}^{J}D_j<Q.
\end{equation}
Indeed, the paired intervals are disjoint and all $C_q$ lie between $-Q/2$ and $Q/2$.

Let $W_0\pmod Q$ be the CRT solution obtained by using the midpoint residue $m_q$ at every unfixed reservoir prime, the selected fixed root at the possible extra reservoir prime, and a fixed root at every $q\in S\setminus S_E$. Choose $W_0$ in $[0,Q)$. For signs
\[
  \sigma,\eta_1,\ldots,\eta_J\in\{-1,1\},
\]
put
\begin{equation}\label{eq:W}
  W_{\sigma,\boldsymbol\eta}
  =W_0+\sigma C_{q_0}+\sum_{j=1}^{J}\eta_jD_j.
\end{equation}
The opposite signs within each adjacent pair ensure that every $W_{\sigma,\boldsymbol\eta}$ satisfies all congruences \eqref{eq:S-congruence}. By \eqref{eq:Cq} and \eqref{eq:width},
\begin{equation}\label{eq:W-range}
  -\frac{3Q}{2}<W_{\sigma,\boldsymbol\eta}<\frac{5Q}{2}.
\end{equation}

For $\alpha\in\mathcal A$, let $Y_{\sigma,\boldsymbol\eta,\alpha}$ be the unique integer satisfying
\begin{equation}\label{eq:Y-CRT}
  0\leq Y_{\sigma,\boldsymbol\eta,\alpha}<GQ,
  \qquad
  Y_{\sigma,\boldsymbol\eta,\alpha}\equiv W_{\sigma,\boldsymbol\eta}\pmod Q,
  \qquad
  Y_{\sigma,\boldsymbol\eta,\alpha}\equiv\alpha\pmod G.
\end{equation}
By \eqref{eq:W-range},
\begin{equation}\label{eq:Y-W-m}
  Y_{\sigma,\boldsymbol\eta,\alpha}
  =W_{\sigma,\boldsymbol\eta}+mQ
\end{equation}
for some
\begin{equation}\label{eq:m-set}
  m\in\{-2,-1,0,1,\ldots,G+1\}.
\end{equation}
The important point is that the set in \eqref{eq:m-set} is fixed, independently of $n$.

\section{Anti-concentration and the controlled insoluble primes}

We use the following direct corollary of the finite-field Hal\'asz inequality of Ferber, Jain, Luh and Samotij \cite[Theorem~1.4]{FJLS}.

\begin{lemma}[Finite-field Hal\'asz bound]\label{lem:Halasz}
Let $p$ be an odd prime, let $a_1,\ldots,a_s$ be nonzero elements of $\F_p$, and let $\epsilon_1,\ldots,\epsilon_s$ be independent uniform signs. Let $R$ be the number of ordered signed relations
\[
  \pm a_i\pm a_j=0
\]
with $i\neq j$. Then
\begin{equation}\label{eq:Halasz}
  \sup_{x\in\F_p}
  \Pr\left(\sum_{i=1}^{s}\epsilon_i a_i=x\right)
  \ll
  \frac1p+\frac1{s^{3/2}}+\frac{R}{s^{5/2}}+e^{-s/100}.
\end{equation}
The implied constant is absolute.
\end{lemma}

\begin{proof}
Apply \cite[Theorem~1.4]{FJLS} with $k=1$, vector length $s$, and $M=s/100$. Its hypotheses are
\[
  30M\leq s,
  \qquad
  80M\leq s,
\]
which hold. The quantity $R_1$ in that theorem counts all ordered signed two-term zero sums, including repetitions. Since $p$ is odd and the $a_i$ are nonzero, the repeated-index relations contribute exactly $2s$, and all remaining relations are counted by $R$. Substitution gives \eqref{eq:Halasz}.
\end{proof}

Fix $p\in T_d$ with $d\leq D$. Since $p\mid P$, put
\[
  z_p\equiv \frac bp
  =N\frac PpQY\pmod p.
\]
The multiplier $N(P/p)Q$ is nonzero modulo $p$. Lemma~\ref{lem:T-local} shows that it is enough to avoid a set
\begin{equation}\label{eq:Bp}
  \mathcal B_p\subseteq\F_p,
  \qquad
  |\mathcal B_p|\leq d\leq D,
\end{equation}
of bad residues for $Y\pmod p$.

We next estimate the concentration of the random sum in \eqref{eq:W}. For a pair $(q_j,r_j)$, equations \eqref{eq:tq} and \eqref{eq:Dj} give
\begin{equation}\label{eq:Dj-factor}
  D_j=\frac{Q}{q_jr_j}
      \bigl(q_jt_{r_j}-r_jt_{q_j}\bigr).
\end{equation}
The integer in parentheses is nonzero: reducing a hypothetical equality modulo $q_j$ would give $q_j\mid t_{q_j}$, contrary to \eqref{eq:tq}. Its absolute value is less than $q_jr_j$.

Every reservoir prime is at most $n/E$, whereas every $p\in T^{(D)}$ is larger than $n/(D+1)$. Since $E>D+1$, a fixed $D_j$ can vanish modulo at most one prime in $T^{(D)}$. Indeed, two such primes would have product larger than $n^2/(D+1)^2$, while the nonzero factor in parentheses in \eqref{eq:Dj-factor} has absolute value less than $n^2/E^2$.

Let
\[
  z(p)=|\{j:D_j\equiv0\pmod p\}|.
\]
Then
\begin{equation}\label{eq:sum-z}
  \sum_{p\in T^{(D)}}z(p)\leq J.
\end{equation}
Consequently, there is at most one prime $p_*$ for which
\begin{equation}\label{eq:exceptional}
  z(p_*)>\frac J2.
\end{equation}
We call it exceptional if it exists.

For a nonexceptional $p$, at least $J/2$ of the coefficients $D_j$ are nonzero modulo $p$. Let $R_p$ be the number of ordered signed relations
\[
  \pm D_j\pm D_k=0\pmod p,
  \qquad j\neq k,
\]
among these nonzero coefficients. Then
\begin{equation}\label{eq:sum-R}
  \sum_{\substack{p\in T^{(D)}\\p\neq p_*}}R_p\ll_{D,E}J^2.
\end{equation}
To prove this, fix two distinct pairs $(q,r)$ and $(u,v)$, where the four letters denote distinct reservoir primes and
\[
  \frac{D_j}{Q}=\frac{t_r}{r}-\frac{t_q}{q},
  \qquad
  \frac{D_k}{Q}=\frac{t_v}{v}-\frac{t_u}{u}.
\]
For signs $\epsilon,\epsilon'\in\{-1,1\}$, a congruence
\[
  \epsilon D_j+\epsilon'D_k\equiv0\pmod p
\]
implies that $p$ divides
\begin{equation}\label{eq:relation-integer}
  F_{\epsilon,\epsilon'}
  =\epsilon uv(qt_r-rt_q)+\epsilon'qr(ut_v-vt_u).
\end{equation}
This integer is nonzero: modulo $q$ it is congruent to
\[
  -\epsilon ruv\,t_q\not\equiv0\pmod q,
\]
because $q,r,u,v$ are distinct primes and $0<|t_q|<q/2$. Moreover,
$|F_{\epsilon,\epsilon'}|=O_E(n^4)$. Since every controlled prime exceeds $n/(D+1)$, a nonzero integer of this size has only $O_D(1)$ controlled prime divisors. Summing over the $O(J^2)$ ordered signed pairs proves \eqref{eq:sum-R}.

Choose the signs $\eta_j$ independently and uniformly. For a nonexceptional $p$, delete the coefficients $D_j$ that vanish modulo $p$ and apply Lemma~\ref{lem:Halasz}. Uniformly in $p$,
\[
  \rho_p
  :=\sup_{x\in\F_p}
       \Pr\left(\sum_{j=1}^{J}\eta_jD_j=x\right)
  \ll
  \frac1p+\frac1{J^{3/2}}+\frac{R_p}{J^{5/2}}+e^{-c_0J}
\]
for an absolute $c_0>0$. Now
\[
  |T^{(D)}|=O_D\!\left(\frac n{\log n}\right)=O_{D,E}(J)
\]
and, since every controlled prime is larger than $n/(D+1)$,
\[
  \sum_{p\in T^{(D)}}\frac1p=O_D\!\left(\frac1{\log n}\right).
\]
Together with \eqref{eq:sum-R}, this gives
\begin{align*}
  \sum_{\substack{p\in T^{(D)}\\p\neq p_*}}\rho_p
  &\ll_{D,E}
    \frac1{\log n}
    +\frac{|T^{(D)}|}{J^{3/2}}
    +\frac{1}{J^{5/2}}
       \sum_{\substack{p\in T^{(D)}\\p\neq p_*}}R_p
    +|T^{(D)}|e^{-c_0J}\\
  &\ll_{D,E}
    \frac1{\log n}+\frac1{\sqrt J}+Je^{-c_0J}
   =o(1).
\end{align*}
Thus
\begin{equation}\label{eq:sum-rho}
  \sum_{\substack{p\in T^{(D)}\\p\neq p_*}}\rho_p=o(1).
\end{equation}

Fix $p\neq p_*$, $\sigma\in\{-1,1\}$, and $\alpha\in\mathcal A$. If the integer $Y_{\sigma,\boldsymbol\eta,\alpha}$ in \eqref{eq:Y-CRT} is bad modulo $p$, then by \eqref{eq:W}, \eqref{eq:Y-W-m}, \eqref{eq:m-set}, and \eqref{eq:Bp}, the random sum
\[
  \sum_{j=1}^{J}\eta_jD_j\pmod p
\]
hits one of at most
\[
  D(G+4)
\]
fixed target residues. Hence the probability of this event is at most $D(G+4)\rho_p$. There are only $2|\mathcal A|=2^{r+1}$ choices of $(\sigma,\alpha)$. By \eqref{eq:sum-rho} and the union bound, for all sufficiently large $n$ there is a choice of $\boldsymbol\eta$ such that
\begin{equation}\label{eq:all-nonexceptional-good}
  Y_{\sigma,\boldsymbol\eta,\alpha}\notin\mathcal B_p
\end{equation}
for every nonexceptional $p\in T^{(D)}$, every $\sigma$, and every $\alpha\in\mathcal A$. Fix such a choice of $\boldsymbol\eta$.

\section{The exceptional prime and the size of \texorpdfstring{$Y$}{Y}}

If no exceptional prime exists, choose any $\alpha_\sigma\in\mathcal A$ for each $\sigma\in\{-1,1\}$. Suppose that $p_*$ exists. For fixed $\sigma$, the residues
\[
  Y_{\sigma,\boldsymbol\eta,\alpha}\pmod{p_*},
  \qquad \alpha\in\mathcal A,
\]
are pairwise distinct. Indeed, two such integers have the same residue modulo $Q$, so their difference is $kQ$ with $|k|<G$. If $p_*$ divides this difference, then $p_*\nmid Q$ and $p_*>G$ for all sufficiently large $n$, so $k=0$. The two integers are then equal, and their residues modulo $G$ are equal, forcing the two values of $\alpha$ to be equal.

By \eqref{eq:A-set} and \eqref{eq:Bp},
\[
  |\mathcal A|=2^r>D\geq|\mathcal B_{p_*}|.
\]
Thus, for each $\sigma$, we may choose $\alpha_\sigma\in\mathcal A$ such that
\[
  Y_\sigma:=Y_{\sigma,\boldsymbol\eta,\alpha_\sigma}
\]
is good modulo $p_*$. By \eqref{eq:all-nonexceptional-good}, both $Y_+$ and $Y_-$ are therefore good at every prime of $T^{(D)}$.

Modulo $Q$,
\[
  Y_+-Y_-\equiv2C_{q_0}.
\]
By \eqref{eq:tq}, every integer in this nonzero residue class has absolute value at least $Q/q_0$: indeed,
\[
  2C_{q_0}+kQ
  =\frac{Q}{q_0}(2t_{q_0}+kq_0),
\]
and the integer in parentheses is nonzero. Hence
\[
  |Y_+-Y_-|\geq\frac{Q}{q_0}\geq\frac Qn.
\]
Both values lie in $[0,GQ)$, so at least one of them satisfies
\begin{equation}\label{eq:Y-size}
  \frac Qn\leq Y<GQ.
\end{equation}
Fix such a $Y$ and define $a,b$ by \eqref{eq:b-form}.

\section{Completion of the proof}

The congruences imposed above imply
\begin{equation}\label{eq:Y-coprime}
  \gcd(Y,NPQ)=1.
\end{equation}
Indeed, the chosen residues modulo every $q\mid Q$ are nonzero by Lemma~\ref{lem:poly}; the chosen residues modulo every $\ell_i$ are also nonzero; and for every $p\mid P$, the residue $Y=0$ corresponds to $z_p=0$, which belongs to the bad set and was avoided.

For every $q\in S_d$, congruence \eqref{eq:S-congruence} says that
\[
  \frac bq=NP\frac QqY
\]
is a root of $f_d$ modulo $q$. Lemma~\ref{lem:S-local} therefore gives
\begin{equation}\label{eq:Q-drop}
  \nu_q(v_{a,b})\leq\nu_q(v_{a,b-1})-1
  \qquad(q\mid Q).
\end{equation}
For every $\ell_i\mid G$, equation \eqref{eq:aux-drop} gives the same decrease. For every $p\mid P$, our choice of $Y$ and Lemma~\ref{lem:T-local} give
\begin{equation}\label{eq:P-neutral}
  \nu_p(v_{a,b})\leq\nu_p(v_{a,b-1}).
\end{equation}
At every other prime, Lemma~\ref{lem:endpoint} and \eqref{eq:Y-coprime} show that a possible increase is at most the exponent of that prime in $Y$. Multiplying the valuation inequalities yields
\begin{equation}\label{eq:ratio}
  \frac{v_{a,b}}{v_{a,b-1}}
  \leq\frac{Y}{GQ}<1
\end{equation}
by \eqref{eq:Y-size}. Hence
\[
  v_{a,b}<v_{a,b-1}.
\]

It remains to estimate the gap. By \eqref{eq:Q-asymp}, \eqref{eq:P-asymp}, \eqref{eq:N-small}, and \eqref{eq:Y-size},
\[
  \log Y=\log Q+O(\log n)=(c+o(1))n
\]
and therefore
\begin{equation}\label{eq:log-b}
\begin{aligned}
  \log b
  &=\log N+\log P+\log Q+\log Y\\
  &=(\tau_D+2c+o(1))n.
\end{aligned}
\end{equation}
In particular, $b$ is exponential in $n$, so
\[
  a=b-n=(1-o(1))b
\]
and
\[
  \log a=(\tau_D+2c+o(1))n.
\]
Since $b-a=n$,
\begin{equation}\label{eq:gap-limit}
  \frac{b-a}{\log a}
  =\frac1{\tau_D+2c}+o(1).
\end{equation}
By \eqref{eq:D-choice}, the right-hand side is less than
\[
  \frac1{1+c}+\varepsilon
\]
for all sufficiently large tuples in Lemma~\ref{lem:simultaneous-powers}. There are infinitely many such tuples, and $a\to\infty$ along them. Since $b(a)\leq b$ for each constructed value of $a$, the asserted upper bound for the lower limit follows as well. This proves Theorem~\ref{thm:main}.

\begin{thebibliography}{9}

\bibitem{VD}
W. van Doorn,
\emph{On the non-monotonicity of the denominator of generalized harmonic sums},
unpublished manuscript.

\bibitem{FJLS}
A. Ferber, V. Jain, K. Luh and W. Samotij,
\emph{On the counting problem in inverse Littlewood--Offord theory},
J. London Math. Soc. (2) \textbf{103} (2021), 1333--1362.
\href{https://doi.org/10.1112/jlms.12409}{doi:10.1112/jlms.12409}.

\bibitem{Chebotarev}
P. Stevenhagen and H. W. Lenstra Jr.,
\emph{Chebotar\"ev and his density theorem},
Math. Intelligencer \textbf{18} (1996), no.~2, 26--37.

\end{thebibliography}

\end{document}
